\sheettitle
  {Self-contained lecture notes}
  {Admissibility, representation, transformation, and finite conditioning}

These notes contain every definition, theorem statement, and proof interface required by the exercises. Self-contained does not mean encyclopedic: the signed Radon--Nikodym theorem is stated in the exact finite regime in which it is used, while its Hahn--Jordan construction remains a declared deferred proof debt.

\section*{1. The probability-space $L^p$ gate}

Let $(E,\mathcal A,\mu)$ be a finite measure space. If $1\le p<q<\infty$, H\"older's inequality with exponents $q/p$ and $q/(q-p)$ gives
\[
\|f\|_p^p
=
\int_E |f|^p\ind_E\,d\mu
\le
\left(\int_E|f|^q\,d\mu\right)^{p/q}
\mu(E)^{1-p/q}.
\]
Hence
\[
\boxed{
\|f\|_p
\le
\mu(E)^{1/p-1/q}\|f\|_q
},
\qquad
L^q(\mu)\subseteq L^p(\mu).
\]
For $q=\infty$, the same conclusion follows from
\[
\int_E|f|^p\,d\mu
\le
\|f\|_\infty^p\mu(E).
\]
On a probability space this becomes
\[
L^\infty(\mathbb P)\subseteq L^q(\mathbb P)\subseteq L^p(\mathbb P)\subseteq L^1(\mathbb P),
\qquad
\|f\|_p\le\|f\|_q
\quad(p<q).
\]
There is no corresponding inclusion on an arbitrary infinite measure space.

Let $Y\in L^1(\mathbb P)$ and let $\mathcal G\subseteq\mathcal F$ be a sub-$\sigma$-field. Define on $(\Omega,\mathcal G)$
\[
\nu_Y^{\mathcal G}(A)
=
\mathbb E[Y\ind_A],
\qquad A\in\mathcal G.
\]
Writing $Y=Y^+-Y^-$ shows that $\nu_Y^{\mathcal G}$ is a finite signed measure. Moreover,
\[
|\nu_Y^{\mathcal G}|(\Omega)
\le
\mathbb E[|Y|]
<\infty,
\qquad
\nu_Y^{\mathcal G}\ll\mathbb P|_{\mathcal G}.
\]

\textbf{Signed Radon--Nikodym interface.}
If $\mu$ is $\sigma$-finite and $\nu$ is a finite signed measure with $\nu\ll\mu$, then there exists $h\in L^1(\mu)$ such that
\[
\nu(A)=\int_Ah\,d\mu
\quad\text{for every measurable }A,
\]
and $h$ is unique $\mu$-almost everywhere. Applied on $(\Omega,\mathcal G)$, this gives a $\mathcal G$-measurable $H\in L^1(\mathbb P|_{\mathcal G})$, unique $\mathbb P$-almost surely, such that
\[
\mathbb E[Y\ind_A]=\int_AH\,d\mathbb P,
\qquad A\in\mathcal G.
\]
This representative is $\mathbb E[Y\mid\mathcal G]$. The numerator, reference measure, measurable space, and almost-everywhere basis are all part of the object.

\section*{2. Expectation and its domain of definition}

For a measurable $X:\Omega\to\overline{\mathbb R}$, write
\[
X^+=\max\{X,0\},
\qquad
X^-=\max\{-X,0\}.
\]
The regimes are distinct:
\begin{itemize}
  \item If $X\ge0$, then $\mathbb E[X]\in[0,+\infty]$ is always defined.
  \item For signed $X$, the expression $\mathbb E[X^+]-\mathbb E[X^-]$ is defined unless both terms are $+\infty$.
  \item If $X\in L^1(\mathbb P)$, then $\mathbb E[X]$ is finite and $|\mathbb E[X]|\le\mathbb E[|X|]$.
\end{itemize}
Thus ``defined'' and ``finite'' are not synonyms. A nonnegative random variable may have expectation $+\infty$; a signed random variable may have no expectation at all because $+\infty-(+\infty)$ is not defined.

\section*{3. Laws and the pushforward integration formula}

Let $X:(\Omega,\mathcal F,\mathbb P)\to(S,\mathcal S)$ be measurable. Its law is the probability measure
\[
\mu_X=X_\#\mathbb P,
\qquad
\mu_X(C)=\mathbb P(X\in C),
\quad C\in\mathcal S.
\]

\textbf{LOTUS.}
For every measurable $\varphi:S\to[0,+\infty]$,
\[
\boxed{
\mathbb E[\varphi(X)]
=
\int_S\varphi\,d\mu_X
}.
\]
The same identity holds for signed $\varphi$ whenever the integral is defined.

\textbf{Proof.}
For $C\in\mathcal S$,
\[
\mathbb E[\ind_C(X)]
=
\mathbb P(X\in C)
=
\mu_X(C)
=
\int_S\ind_C\,d\mu_X.
\]
Linearity gives the identity for nonnegative simple functions. If $\varphi\ge0$ is measurable, choose nonnegative simple functions $\varphi_n\uparrow\varphi$. Then $\varphi_n(X)\uparrow\varphi(X)$, so monotone convergence on both spaces yields
\[
\mathbb E[\varphi(X)]
=
\lim_n\mathbb E[\varphi_n(X)]
=
\lim_n\int_S\varphi_n\,d\mu_X
=
\int_S\varphi\,d\mu_X.
\]
For signed $\varphi$, apply the nonnegative result to $\varphi^+$ and $\varphi^-$ whenever their difference is defined. \hfill$\square$

If $\mu_X$ is atomic, the integral is a weighted sum. If $\mu_X\ll\lambda$ on $\mathbb R^d$, Radon--Nikodym supplies a density $f_X=d\mu_X/d\lambda$, unique $\lambda$-almost everywhere, and
\[
\mathbb E[\varphi(X)]
=
\int_{\mathbb R^d}\varphi(x)f_X(x)\,dx.
\]
The density formula is a representation of LOTUS under an additional domination hypothesis. The law itself needs no chosen reference measure and no density.

\section*{4. Tonelli and the layer-cake representation}

Let $(E,\mathcal A,\mu)$ and $(F,\mathcal B,\nu)$ be $\sigma$-finite measure spaces. If $f:E\times F\to[0,+\infty]$ is measurable, Tonelli's theorem gives
\[
\int_{E\times F}f\,d(\mu\otimes\nu)
=
\int_E\!\left(\int_Ff(x,y)\,\nu(dy)\right)\mu(dx)
=
\int_F\!\left(\int_Ef(x,y)\,\mu(dx)\right)\nu(dy),
\]
with all three values allowed to be $+\infty$. Positivity is decisive. For a signed integrand, Fubini requires absolute integrability; otherwise an exchange may conceal the undefined expression $+\infty-(+\infty)$.

Let $f:E\to[0,+\infty]$ be measurable and set
\[
H_f=\{(x,t)\in E\times(0,+\infty):t<f(x)\}.
\]
The subgraph is measurable because
\[
H_f
=
\bigcup_{q\in\mathbb Q_{>0}}
\{x:f(x)>q\}\times(0,q).
\]
Since
\[
f(x)=\int_0^\infty\ind_{\{t<f(x)\}}\,dt,
\]
Tonelli applied to $\ind_{H_f}$ gives
\[
\boxed{
\int_Ef\,d\mu
=
\int_0^\infty\mu(f>t)\,dt
}.
\]
More generally, for $p>0$,
\[
\int_Ef^p\,d\mu
=
p\int_0^\infty t^{p-1}\mu(f>t)\,dt.
\]
On a probability space, for $X\ge0$,
\[
\boxed{
\mathbb E[X]
=
\int_0^\infty\mathbb P(X>t)\,dt
}.
\]
The identity remains valid when the common value is $+\infty$.

\section*{5. Transformed insurance payments}

For $X\ge0$, a deductible $d>0$, and a cap $u>0$, apply the tail formula to the transformed variable, not to its mean. Since
\[
\mathbb P((X-d)_+>t)=\mathbb P(X>d+t),
\]
one obtains
\[
\boxed{
\mathbb E[(X-d)_+]
=
\int_d^\infty\mathbb P(X>s)\,ds
}.
\]
Similarly,
\[
\boxed{
\mathbb E[X\wedge u]
=
\int_0^u\mathbb P(X>s)\,ds
},
\]
and the limited-loss payment satisfies
\[
\boxed{
\mathbb E[\min\{(X-d)_+,u\}]
=
\int_d^{d+u}\mathbb P(X>s)\,ds
}.
\]
If $X\sim\operatorname{Exp}(\lambda)$, these become
\[
\mathbb E[(X-d)_+]
=
\frac{e^{-\lambda d}}{\lambda},
\qquad
\mathbb E[X\wedge u]
=
\frac{1-e^{-\lambda u}}{\lambda},
\]
and
\[
\mathbb E[\min\{(X-d)_+,u\}]
=
\frac{e^{-\lambda d}(1-e^{-\lambda u})}{\lambda}.
\]

If $g$ is convex, $X\in L^1$, and $g(X)\in L^1$, Jensen's inequality gives
\[
g(\mathbb E[X])\le\mathbb E[g(X)].
\]
It gives a comparison, not the value of $\mathbb E[g(X)]$.

\section*{6. Conditioning on a positive-probability event}

For $B\in\mathcal F$ with $\mathbb P(B)>0$, define the normalized restriction
\[
\mathbb P_B(A)
=
\frac{\mathbb P(A\cap B)}{\mathbb P(B)},
\qquad A\in\mathcal F.
\]
Then $\mathbb P_B\ll\mathbb P$ and
\[
\boxed{
\frac{d\mathbb P_B}{d\mathbb P}
=
\frac{\ind_B}{\mathbb P(B)}
}
\quad\mathbb P\text{-almost surely}.
\]
For $Y\in L^1(\mathbb P)$,
\[
\mathbb E[Y\mid B]
=
\frac{\mathbb E[Y\ind_B]}{\mathbb P(B)}.
\]
The ratio formula does not condition on a null event. If $X$ has a continuous law, then typically $\mathbb P(X=x)=0$; conditioning on $X=x$ requires a conditional kernel or density and a version choice, which lies beyond this packet.

\section*{7. Finite partitions, total expectation, and Bayes}

Let $(B_i)_{1\le i\le r}$ be a measurable partition and set $\mathcal G=\sigma(B_1,\ldots,B_r)$. Every finite real-valued $\mathcal G$-measurable random variable is constant on each atom, hence has the form
\[
H=\sum_{i=1}^rc_i\ind_{B_i}.
\]
For $Y\in L^1(\mathbb P)$, the conditional-expectation identities
\[
\int_AH\,d\mathbb P
=
\int_AY\,d\mathbb P,
\qquad A\in\mathcal G,
\]
determine the coefficients on positive-probability atoms:
\[
\boxed{
\mathbb E[Y\mid\mathcal G]
=
\sum_{i:\,\mathbb P(B_i)>0}
\frac{\mathbb E[Y\ind_{B_i}]}{\mathbb P(B_i)}\ind_{B_i}
}.
\]
On a null atom the coefficient is arbitrary; conditional expectation is an almost-sure equivalence class.

Taking expectations gives the law of total expectation,
\[
\mathbb E[Y]
=
\sum_i\mathbb P(B_i)\mathbb E[Y\mid B_i],
\]
where null-atom terms contribute zero. Setting $Y=\ind_A$ gives total probability,
\[
\mathbb P(A)
=
\sum_i\mathbb P(A\mid B_i)\mathbb P(B_i).
\]
If $\mathbb P(A)>0$, normalization of the joint masses gives Bayes' formula:
\[
\boxed{
\mathbb P(B_j\mid A)
=
\frac{\mathbb P(A\mid B_j)\mathbb P(B_j)}
{\sum_i\mathbb P(A\mid B_i)\mathbb P(B_i)}
}.
\]
Bayes and total expectation are consequences of the same finite-partition structure, not additional axioms.

\section*{8. Representation ledger}

A formula is not enough to identify a density. Record its numerator, reference measure, domain, and uniqueness basis. A pushforward law is a measure in its own right and therefore occupies a representation row, not a malformed density row.

{\footnotesize
\hbadness=10000
\begin{tabularx}{\textwidth}{@{}Y Y Y Y Y@{}}
\toprule
Object & Numerator or construction & Reference & Domain & Equality basis \\
\midrule
$d(f\mathbb P)/d\mathbb P=f$ & $A\mapsto\mathbb E[f\ind_A]$ & $\mathbb P$ & $(\Omega,\mathcal F)$ & $\mathbb P$-a.s. \\
$\mu_X=X_\#\mathbb P$ & pushforward of $\mathbb P$ by $X$ & none & $(S,\mathcal S)$ & equality of measures \\
$f_X=d\mu_X/d\lambda$ & law $\mu_X$, only if $\mu_X\ll\lambda$ & $\lambda$ & state space & $\lambda$-a.e. \\
$d\mathbb P_A/d\mathbb P$ & $C\mapsto\mathbb P(C\cap A)/\mathbb P(A)$ & $\mathbb P$ & $(\Omega,\mathcal F)$ & $\mathbb P$-a.s. \\
$\mathbb E[Y\mid\mathcal G]$ & $C\mapsto\mathbb E[Y\ind_C]$, $C\in\mathcal G$ & $\mathbb P|_{\mathcal G}$ & $(\Omega,\mathcal G)$ & $\mathbb P$-a.s. \\
\bottomrule
\end{tabularx}
\par
}

The same algebraic expression may represent different objects on different measurable spaces. The ambient measure and almost-everywhere relation are not decorative annotations; they are part of the statement.
